\section{Product rules (generalized distributive laws)}

\begin{proposition}
\label{prop.fps.prodrule-fin-fin}
\lean{prod_sum_eq_sum_prod_finset}
\leantarget
\leanok
Let $L$ be a commutative ring. For every
$n\in\mathbb{N}$, let $\left[  n\right]  $ denote the set $\left\{
1,2,\ldots,n\right\}  $.

Let $n\in\mathbb{N}$. For every $i\in\left[  n\right]  $, let $p_{i,1}%
,p_{i,2},\ldots,p_{i,m_{i}}$ be finitely many elements of $L$. Then,%
\begin{equation}
\prod_{i=1}^{n}\ \ \sum_{k=1}^{m_{i}}p_{i,k}=\sum_{\left(  k_{1},k_{2}%
,\ldots,k_{n}\right)  \in\left[  m_{1}\right]  \times\left[  m_{2}\right]
\times\cdots\times\left[  m_{n}\right]  }\ \ \prod_{i=1}^{n}p_{i,k_{i}}.
\label{eq.lem.fps.prodrule-fin-fin.eq}%
\end{equation}
\end{proposition}

\begin{proof}
The idea is to reduce it to the
case $n=2$ by induction, then to use the discrete Fubini rule.
\end{proof}

\begin{proposition}
\label{prop.fps.prodrule-fin-inf}
\lean{prod_tsum_eq_tsum_prod_finset}
\leantarget
\leanok
For every $n\in\mathbb{N}$, let $\left[
n\right]  $ denote the set $\left\{  1,2,\ldots,n\right\}  $.

Let $n\in\mathbb{N}$. For every $i\in\left[  n\right]  $, let $\left(
p_{i,k}\right)  _{k\in S_{i}}$ be a summable family of elements of $K\left[
\left[  x\right]  \right]  $. Then,%
\begin{equation}
\prod_{i=1}^{n}\ \ \sum_{k\in S_{i}}p_{i,k}=\sum_{\left(  k_{1},k_{2}%
,\ldots,k_{n}\right)  \in S_{1}\times S_{2}\times\cdots\times S_{n}}%
\ \ \prod_{i=1}^{n}p_{i,k_{i}}. \label{eq.lem.fps.prodrule-fin-inf.eq}%
\end{equation}
In particular, the family $\left(  \prod_{i=1}^{n}p_{i,k_{i}}\right)
_{\left(  k_{1},k_{2},\ldots,k_{n}\right)  \in S_{1}\times S_{2}\times
\cdots\times S_{n}}$ is summable.
\end{proposition}

\begin{proof}
Same method as for Proposition \ref{prop.fps.prodrule-fin-fin}, but now using
the discrete Fubini rule for infinite sums.
The proof is by induction on $n$:
\begin{itemize}
\item Base case ($n=0$): Both sides equal $1$.
\item Inductive step: Split the product into the first $n$ factors and the last factor,
apply the induction hypothesis on the first $n$, then apply the Cauchy
product formula for infinite sums.
\end{itemize}
Summability of the product family is established by showing that for each
coefficient $d$, only finitely many functions $f$ give a nonzero contribution
(using finite support of each summable family in the discrete topology).
\end{proof}

\subsubsection{Helpers for the finite product rule}

\begin{lemma}
\label{lem.fps.coeff-finProd-zero-of-high-order}
\lean{coeff_finProd_eq_zero_of_factor_high_order}
\leanhelper
\leanok
Let $n\in\mathbb{N}$. Let $\left(p_{i,k}\right)$ be a family indexed by
$\{1,2,\ldots,n\}\times S_i$, and let $f$ be a choice function selecting
$f(i) \in S_i$ for each $i$.
If for some $j\in\{1,2,\ldots,n\}$, the FPS $p_{j,f(j)}$ satisfies
$[x^m](p_{j,f(j)}) = 0$ for all $m\le d$, then
$[x^d]\left(\prod_{i} p_{i,f(i)}\right) = 0$.
\end{lemma}

\begin{proof}
The product $\prod_i p_{i,f(i)}$ is divisible by $p_{j,f(j)}$, which has
order $> d$. Since the order of a product is at least the sum of the orders
of the factors, the product also has order $> d$.
\end{proof}

\begin{lemma}
\label{lem.fps.summable-finProd-discrete}
\lean{summable_finProd_discrete}
\leanhelper
\leanok
\textbf{(Summability of product families.)}
Assume $K$ has the discrete topology. Let $n\in\mathbb{N}$,
and for each $i\in\{1,2,\ldots,n\}$, let $\left(p_{i,k}\right)_{k\in S_i}$
be a summable family in $K[[x]]$.
Then the family $\left(\prod_i p_{i,f(i)}\right)_{f\in\prod_i S_i}$
is summable.
\end{lemma}

\begin{proof}
For each coefficient $d$, define $T_i = \bigcup_{m\le d}
\mathrm{supp}([x^m]\circ p_i)$. Each $T_i$ is finite (by discrete summability).
If $f(i)\notin T_i$ for some $i$, then $[x^d](\prod_j p_{j,f(j)}) = 0$
by Lemma \ref{lem.fps.coeff-finProd-zero-of-high-order}.
The set $\{f \mid \forall i,\, f(i)\in T_i\}$ is finite (product of finite
sets), so the support at coefficient $d$ is finite.
\end{proof}

\begin{lemma}
\label{lem.fps.summable-prod-of-summable-discrete}
\lean{summable_prod_of_summable_discrete}
\leanhelper
\leanok
\textbf{(Product of two summable families.)}
Assume $K$ has the discrete topology. If $(f_\alpha)_{\alpha\in A}$ and
$(g_\beta)_{\beta\in B}$ are summable families in $K[[x]]$, then the family
$(f_\alpha\cdot g_\beta)_{(\alpha,\beta)\in A\times B}$ is summable.
\end{lemma}

\begin{proof}
\leanok
For each coefficient $d$, the support of
$(\alpha,\beta)\mapsto [x^d](f_\alpha g_\beta)$ is contained in the
finite product of coefficient supports, using the Cauchy product formula
and the finiteness of supports.
\end{proof}

\begin{definition}
\label{def.fps.prodrule.ess-fin}
\lean{EssentiallyFinite}
\leantarget
\leanok
\textbf{(a)} A sequence $\left(  k_{1}%
,k_{2},k_{3},\ldots\right)  $ is said to be \emph{essentially finite} if all
but finitely many $i\in\left\{  1,2,3,\ldots\right\}  $ satisfy $k_{i}=0$.
\medskip

\textbf{(b)} A family $\left(  k_{i}\right)  _{i\in I}$ is said to be
\emph{essentially finite} if all but finitely many $i\in I$ satisfy $k_{i}=0$.
\end{definition}

\begin{proposition}
\label{prop.fps.prodrule-inf-infN}
\lean{tprod_tsum_eq_tsum_prod_essentiallyFinite_nat}
\leantarget
\leanok
Let $S_{1},S_{2},S_{3},\ldots$ be infinitely
many sets that all contain the number $0$. Set%
\[
\overline{S}=\left\{  \left(  i,k\right)  \ \mid\ i\in\left\{  1,2,3,\ldots
\right\}  \text{ and }k\in S_{i}\text{ and }k\neq0\right\}  .
\]

For any $i\in\left\{  1,2,3,\ldots\right\}  $ and any $k\in S_{i}$, let
$p_{i,k}$ be an element of $K\left[  \left[  x\right]  \right]  $. Assume
that
\begin{equation}
p_{i,0}=1\ \ \ \ \ \ \ \ \ \ \text{for any }i\in\left\{  1,2,3,\ldots\right\}
. \label{eq.prop.fps.prodrule-inf-infN.pi0=1}%
\end{equation}
Assume further that the family $\left(  p_{i,k}\right)  _{\left(  i,k\right)
\in\overline{S}}$ is summable. Then, the product $\prod_{i=1}^{\infty}%
\ \ \sum_{k\in S_{i}}p_{i,k}$ is well-defined (i.e., the family $\left(
p_{i,k}\right)  _{k\in S_{i}}$ is summable for each $i\in\left\{
1,2,3,\ldots\right\}  $, and the family $\left(  \sum_{k\in S_{i}}%
p_{i,k}\right)  _{i\in\left\{  1,2,3,\ldots\right\}  }$ is multipliable), and
we have%
\begin{equation}
\prod_{i=1}^{\infty}\ \ \sum_{k\in S_{i}}p_{i,k}=\sum_{\substack{\left(
k_{1},k_{2},k_{3},\ldots\right)  \in S_{1}\times S_{2}\times S_{3}\times
\cdots\\\text{is essentially finite}}}\ \ \prod_{i=1}^{\infty}p_{i,k_{i}}.
\label{eq.prop.fps.prodrule-inf-infN.eq}%
\end{equation}
In particular, the family $\left(  \prod_{i=1}^{\infty}p_{i,k_{i}}\right)
_{\left(  k_{1},k_{2},k_{3},\ldots\right)  \in S_{1}\times S_{2}\times
S_{3}\times\cdots\text{ is essentially finite}}$ is summable.
\end{proposition}

\begin{proof}
This is the particular case of Proposition \ref{prop.fps.prodrule-inf-inf}
for $I=\left\{  1,2,3,\ldots\right\}  $.
\end{proof}

\begin{proposition}
\label{prop.fps.prodrule-inf-inf}
\lean{tprod_tsum_eq_tsum_prod_essentiallyFinite}
\leantarget
\leanok
Let $I$ be a set. For any $i\in I$, let
$S_{i}$ be a set that contains the number $0$. Set%
\[
\overline{S}=\left\{  \left(  i,k\right)  \ \mid\ i\in I\text{ and }k\in
S_{i}\text{ and }k\neq0\right\}  .
\]

For any $i\in I$ and any $k\in S_{i}$, let $p_{i,k}$ be an element of
$K\left[  \left[  x\right]  \right]  $. Assume that
\begin{equation}
p_{i,0}=1\ \ \ \ \ \ \ \ \ \ \text{for any }i\in I.
\label{eq.prop.fps.prodrule-inf-inf.pi0=1}%
\end{equation}
Assume further that the family $\left(  p_{i,k}\right)  _{\left(  i,k\right)
\in\overline{S}}$ is summable. Then, the product $\prod_{i\in I}\ \ \sum_{k\in
S_{i}}p_{i,k}$ is well-defined (i.e., the family $\left(  p_{i,k}\right)
_{k\in S_{i}}$ is summable for each $i\in I$, and the family $\left(
\sum_{k\in S_{i}}p_{i,k}\right)  _{i\in I}$ is multipliable), and we have%
\begin{equation}
\prod_{i\in I}\ \ \sum_{k\in S_{i}}p_{i,k}=\sum_{\substack{\left(
k_{i}\right)  _{i\in I}\in\prod_{i\in I}S_{i}\\\text{is essentially finite}%
}}\ \ \prod_{i\in I}p_{i,k_{i}}. \label{eq.prop.fps.prodrule-inf-inf.eq}%
\end{equation}
In particular, the family $\left(  \prod_{i\in I}p_{i,k_{i}}\right)  _{\left(
k_{i}\right)  _{i\in I}\in\prod_{i\in I}S_{i}\text{ is essentially finite}}$
is summable.
\end{proposition}

\begin{proof}
The proof is a long-winded reduction to the finite case.
It proceeds by showing both sides have the same $n$-th coefficient for each $n$.

\textbf{Setup.} For each $m\in\mathbb{N}$, define
$T'_n = \{(i,k)\in\overline{S} \mid \exists m\le n,\, [x^m](p_{i,k})\ne 0\}$.
This set is finite by coefficient-wise summability.
Let $I_n = \{i \in I \mid \exists k: (i,k)\in T'_n\}$.
Then $I_n$ is finite.

\textbf{RHS reduction (Claim 8).}
For the RHS, only essentially finite functions $f$ with ``graph'' contained in $T'_n$
can contribute to the $n$-th coefficient. If some $(j, f(j))\notin T'_n$ with $f(j)\ne 0$,
then $[x^m](p_{j,f(j)}) = 0$ for all $m\le n$, so $[x^n](\prod_i p_{i,f(i)}) = 0$
(by the finite product coefficient lemma).
The set of such functions is finite
(since their graphs inject into the finite powerset of $T'_n$).

\textbf{LHS reduction (Claim 10).}
For the LHS, write $\sum_k p_{i,k} = 1 + \sum_{k\ne 0} p_{i,k}$.
For $i\notin I_n$, the sum $\sum_{k\ne 0} p_{i,k}$ has all coefficients
up to $n$ equal to $0$. By the irrelevance lemma (Lemma \ref{lem.fps.prod.irlv.inf}),
the factors for $i\notin I_n$ do not affect the $n$-th coefficient.
Thus $[x^n](\prod_i \sum_k p_{i,k}) = [x^n](\prod_{i\in I_n} \sum_k p_{i,k})$.

\textbf{Finite product rule (Claim 11).}
Since $I_n$ is finite, we can apply Proposition \ref{prop.fps.prodrule-fin-infJ}:
$\prod_{i\in I_n} \sum_k p_{i,k} = \sum_{(k_i)_{i\in I_n}\in\prod_{i\in I_n} S_i}
\prod_{i\in I_n} p_{i,k_i}$.

Combining Claims 8, 10, and 11 shows both sides have equal $n$-th coefficients.
\end{proof}

\begin{proposition}
\label{prop.fps.prodrule-fin-infJ}
\lean{prod_tsum_eq_tsum_prod_finset'_discrete}
\leantarget
\leanok
Let $N$ be a finite set. For any $i\in N$,
let $\left(  p_{i,k}\right)  _{k\in S_{i}}$ be a summable family of elements
of $K\left[  \left[  x\right]  \right]  $. Then,%
\begin{equation}
\prod_{i\in N}\ \ \sum_{k\in S_{i}}p_{i,k}=\sum_{\left(  k_{i}\right)  _{i\in
N}\in\prod_{i\in N}S_{i}}\ \ \prod_{i\in N}p_{i,k_{i}}.
\label{eq.lem.fps.prodrule-fin-infJ.eq}%
\end{equation}
In particular, the family $\left(  \prod_{i\in N}p_{i,k_{i}}\right)  _{\left(
k_{i}\right)  _{i\in N}\in\prod_{i\in N}S_{i}}$ is summable.
\end{proposition}

\begin{proof}
This is just
Proposition \ref{prop.fps.prodrule-fin-inf}, with the indexing set $\left\{
1,2,\ldots,n\right\}  $ replaced by $N$. It can be proved by reindexing the
products (using a bijection between $N$ and $\{1,2,\ldots,|N|\}$) or directly by induction on $\left\vert N\right\vert $.
\end{proof}

\begin{lemma}
\label{lem.fps.prod.irlv.inf}
\lean{coeff_mul_tprod_one_add_eq_coeff}
\leantarget
\leanok
Let $a\in K\left[  \left[  x\right]  \right]  $
be an FPS. Let $\left(  f_{i}\right)  _{i\in J}\in K\left[  \left[  x\right]
\right]  ^{J}$ be a summable family of FPSs. Let $n\in\mathbb{N}$. Assume that
each $i\in J$ satisfies
\begin{equation}
\left[  x^{m}\right]  \left(  f_{i}\right)  =0\ \ \ \ \ \ \ \ \ \ \text{for
each }m\in\left\{  0,1,\ldots,n\right\}  .
\label{eq.lem.fps.prod.irlv.inf.ass}%
\end{equation}
Then,
\[
\left[  x^{m}\right]  \left(  a\prod_{i\in J}\left(  1+f_{i}\right)  \right)
=\left[  x^{m}\right]  a\ \ \ \ \ \ \ \ \ \ \text{for each }m\in\left\{
0,1,\ldots,n\right\}  .
\]
\end{lemma}

\begin{proof}
The idea is to argue that the first $n+1$ coefficients of $a\prod_{i\in J}\left(
1+f_{i}\right)  $ agree with those of $a$.

First, one shows that for any finite subset $s$ of $J$, if each $f_i$ has
$[x^m](f_i) = 0$ for all $m\le n$, then
$[x^m](\prod_{i\in s}(1+f_i)) = [x^m](1)$ for all $m\le n$.
This uses the expansion $\prod_{i\in s}(1+f_i) = \sum_{t\subseteq s}\prod_{i\in t}f_i$;
for $t\ne\emptyset$, each $\prod_{i\in t}f_i$ has order $>n$, so its $m$-th
coefficient is zero.

Then, for the infinite product, the result follows by taking limits:
$\prod_{i\in J}(1+f_i)$ is the limit of the partial products $\prod_{i\in s}(1+f_i)$,
and since $[x^m]$ is continuous, each partial product gives
$[x^m](\prod_{i\in s}(1+f_i)) = [x^m](1)$.

Finally, the coefficient of $a\cdot\prod_{i\in J}(1+f_i)$ uses the Cauchy product formula,
which only involves coefficients of $\prod_{i\in J}(1+f_i)$ up to degree $m\le n$,
all of which equal the corresponding coefficients of $1$.
Hence $[x^m](a\cdot\prod_{i\in J}(1+f_i)) = [x^m](a\cdot 1) = [x^m](a)$.

If the family $(1+f_i)_{i\in J}$ is not multipliable,
the infinite product defaults to $1$ by convention, and the result holds trivially.
\end{proof}

\subsection{Helper lemmas for the infinite product rule}

These lemmas implement the key technical steps in the proof of Proposition
\ref{prop.fps.prodrule-inf-inf}. They require $K$ to have the discrete topology.

\begin{lemma}
\label{lem.fps.T'n-finite-discrete}
\lean{T'n_finite_of_discrete}
\leanhelper
\leanok
\textbf{(Finiteness of $T'_n$.)}
Under the assumptions of Proposition \ref{prop.fps.prodrule-inf-inf}
(with discrete $K$), define
$T'_n = \{(i,k)\in\overline{S} \mid \exists m\le n,\, [x^m](p_{i,k})\ne 0\}$.
Then $T'_n$ is a finite set.
\end{lemma}

\begin{proof}
We have $T'_n = \bigcup_{m\le n}
\mathrm{supp}\!\left((i,k)\mapsto [x^m](p_{i,k})\right)$.
Each support set is finite by coefficient-wise summability
(summable families over a discrete topological space have finite support).
A finite union of finite sets is finite.
\end{proof}

\begin{lemma}
\label{lem.fps.coeff-prod-zero-of-factor-coeff-zero}
\lean{coeff_prod_eq_zero_of_factor_coeff_zero}
\leanhelper
\leanok
If $j\in S$ and $[x^m](f_j) = 0$ for all $m\le n$, then
$[x^m]\!\left(\prod_{i\in S} f_i\right) = 0$ for all $m\le n$.
\end{lemma}

\begin{proof}
\leanok
Write $\prod_{i\in S} f_i = f_j \cdot \prod_{i\in S\setminus\{j\}} f_i$.
For each $m\le n$, use the Cauchy product formula: each term
involves $[x^a](f_j)$ with $a\le m\le n$, which is $0$ by hypothesis.
\end{proof}

\begin{lemma}
\label{lem.fps.finite-funcs-graph-in-finite-set}
\lean{finite_funcs_with_graph_in_finite_set}
\leanhelper
\leanok
Let $T$ be a finite set of pairs $(i,k)$ with $k\ne 0$.
Then the set of essentially finite functions $f : I\to\prod_i S_i$
whose ``graph'' $\{(i,f(i)) : f(i)\ne 0\}$ is contained in $T$
is finite.
\end{lemma}

\begin{proof}
\leanok
Each such function is determined by its graph, which is a subset of $T$.
The map sending $f$ to its graph is injective (since the graph determines $f$
on nonzero values, and the rest are $0$). Since the powerset of $T$ is finite,
the image is finite, so the set of such functions is finite.
\end{proof}

\begin{lemma}
\label{lem.fps.coeff-prod-extra-high-order}
\lean{coeff_prod_eq_of_extra_high_order'}
\leanhelper
\leanok
\textbf{(Finite irrelevance for extra high-order factors.)}
Let $s\subseteq t$ be finite sets of indices, and let $g : I\to K[[x]]$.
Assume that for each $i\in t\setminus s$, we have
$[x^m](g_i) = 0$ for all $m\le n$.
Then $[x^n]\!\left(\prod_{i\in t}(1+g_i)\right) =
[x^n]\!\left(\prod_{i\in s}(1+g_i)\right)$.
\end{lemma}

\begin{proof}
Write $\prod_{i\in t}(1+g_i) = \prod_{i\in s}(1+g_i) \cdot
\prod_{i\in t\setminus s}(1+g_i)$. For the second factor, each $g_i$
($i\in t\setminus s$) has order $>n$, so by expanding
$\prod_{i\in t\setminus s}(1+g_i)$ via the powerset formula, all terms
except the empty product have order $>n$.
Hence $[x^m]\!\left(\prod_{i\in t\setminus s}(1+g_i)\right) =
[x^m](1)$ for all $m\le n$, and the Cauchy product formula gives the result.
\end{proof}

\begin{lemma}
\label{lem.fps.coeff-tprod-one-add-reduce-to-finite}
\lean{coeff_tprod_one_add_eq_coeff_prod_of_high_order'}
\leanhelper
\leanok
\textbf{(Infinite irrelevance: reduction to finite product.)}
Assume $K$ has the discrete topology. Let $(g_i)_{i\in I}$ be a family such
that $(1+g_i)_{i\in I}$ is multipliable, and let $I_n\subseteq I$ be a finite
set such that $[x^m](g_i) = 0$ for all $i\notin I_n$ and all $m\le n$.
Then
\[
[x^n]\!\left(\prod_{i\in I}(1+g_i)\right) =
[x^n]\!\left(\prod_{i\in I_n}(1+g_i)\right).
\]
\end{lemma}

\begin{proof}
\leanok
The infinite product $\prod_{i\in I}(1+g_i)$ is the limit of partial products
$\prod_{i\in s}(1+g_i)$ over finite sets $s$.
By Lemma \ref{lem.fps.coeff-prod-extra-high-order}, for any $s\supseteq I_n$,
the $n$-th coefficient of $\prod_{i\in s}(1+g_i)$ equals that of
$\prod_{i\in I_n}(1+g_i)$.
By continuity of $[x^n]$ and the eventual constancy, the limit also has this
$n$-th coefficient.
\end{proof}

\subsection{Helper lemmas from the Lean formalization}

\begin{lemma}
\label{lem.summable-fiber-discrete}
\lean{summable_fiber_discrete}
\leanhelper
\leanok
\textbf{(Summability of fibers.)}
Let $I$ be a set. For any $i\in I$, let $S_i$ be a set containing $0$.
For any $i\in I$ and $k\in S_i$, let $p_{i,k}\in K[[x]]$.
Assume that the family $(p_{i,k})_{(i,k)\in\overline{S}}$ is summable
(where $\overline{S} = \{(i,k) : k\ne 0\}$).
Then, for each $i\in I$, the family $(p_{i,k})_{k\in S_i}$ is summable.
\end{lemma}

\begin{proof}
The subfamily $(p_{i,k})_{k\in S_i, k\ne 0}$ is summable because it injects
into the summable family $(p_{i,k})_{(i,k)\in\overline{S}}$ via
$(i,\cdot)$. Adding the single term $p_{i,0}$ preserves summability.
Formally, we reduce to coefficient-wise summability and split each coefficient
sum at $k=0$.
\end{proof}

\begin{lemma}
\label{lem.multipliable-tsum-fiber-discrete}
\lean{multipliable_tsum_fiber_discrete}
\leanhelper
\leanok
\textbf{(Multipliability of the sum family.)}
Under the same assumptions as Lemma \ref{lem.summable-fiber-discrete},
and assuming $p_{i,0}=1$ for all $i\in I$,
the family $\left(\sum_{k\in S_i} p_{i,k}\right)_{i\in I}$ is multipliable.
\end{lemma}

\begin{proof}
\leanok
Write $\sum_k p_{i,k} = 1 + \sum_{k\ne 0} p_{i,k}$.
We need to show $(1 + g_i)_{i\in I}$ is multipliable,
where $g_i = \sum_{k\ne 0} p_{i,k}$.
By the multipliability criterion (Theorem \ref{thm.fps.1+f-mulable}), it suffices to show
that $\sum_{s \in \mathcal{F}(I)} \prod_{i\in s} g_i$ is summable,
where $\mathcal{F}(I)$ denotes the set of all finite subsets of $I$.
This is proved coefficient-wise: for each $n$, the support of
$s\mapsto [x^n](\prod_{i\in s}g_i)$ is contained in the set of finite subsets
$s$ with $s\subseteq I_n$ (the finite set of indices contributing to
coefficients up to $n$), which is finite.
\end{proof}

\begin{lemma}
\label{lem.summable-prod-essentiallyFinite-discrete}
\lean{summable_prod_essentiallyFinite_discrete}
\leanhelper
\leanok
\textbf{(Summability of the RHS.)}
Under the assumptions of Proposition \ref{prop.fps.prodrule-inf-inf} (with discrete $K$),
the family $\left(\prod_{i\in I} p_{i,k_i}\right)_{(k_i)\text{ ess.\ fin.}}$
is summable.
\end{lemma}

\begin{proof}
\leanok
For each coefficient $n$, define
$T'_n = \{(i,k)\in\overline{S} \mid \exists m\le n,\, [x^m](p_{i,k})\ne 0\}$.
This is finite by Lemma \ref{lem.fps.T'n-finite-discrete}.
The support of $f\mapsto [x^n](\prod_{i\in I} p_{i,f(i)})$ is contained in
the set of essentially finite functions whose ``graph''
$\{(i,f(i)) : f(i)\ne 0\}$ is a subset of $T'_n$
(if $(j,f(j))\notin T'_n$ for some $j$ with $f(j)\ne 0$, then
$[x^n](\prod_{i\in I} p_{i,f(i)}) = 0$ by Lemma
\ref{lem.fps.coeff-prod-zero-of-factor-coeff-zero}).
This set of functions is finite by Lemma
\ref{lem.fps.finite-funcs-graph-in-finite-set}.
\end{proof}

\subsection{Binary product rule}

\begin{proposition}
\label{prop.fps.binary-product-rule}
\lean{binary_product_rule}
\leanhelper
\leanok
Assume $K$ has the discrete topology. Let $\left(a_i\right)_{i\in\mathbb{N}}$
be a summable family in $K[[x]]$. Then
\[
\prod_{i\in\mathbb{N}}\left(1+a_i\right) =
\sum_{S\subseteq\mathbb{N},\;S\text{ finite}} \prod_{i\in S} a_i.
\]
In particular, the right-hand side is summable.
\end{proposition}

\begin{proof}
\leanok
The identity $\prod_{i\in\mathbb{N}}(1+f_i) = \sum_{s\subseteq\mathbb{N},\,s\text{ finite}} \prod_{i\in s} f_i$ holds
provided the right-hand side is summable.
For the summability: for each coefficient $d$, define
$T = \bigcup_{m\le d}\mathrm{supp}([x^m]\circ a)$, which is finite by
discrete summability. Any finset $s$ not contained in $T$ contributes
$0$ to coefficient $d$ (since some factor $a_i$ with $i\notin T$ has
all low coefficients zero). The set of finsets contained in $T$ is finite.
\end{proof}

\begin{proposition}
\label{prop.fps.binary-product-rule-alt}
\lean{binary_product_rule'}
\leanhelper
\leanok
Alternative formulation: under the same assumptions,
\[
\prod_{i\in\mathbb{N}}\left(1+a_i\right) =
\sum_{\substack{J\subseteq\mathbb{N},\; J\text{ finite}}}
\prod_{i\in J}^{\mathrm{fin}} a_i,
\]
where the RHS sums over all finite subsets $J$ of $\mathbb{N}$
and uses the finitary product $\prod^{\mathrm{fin}}$.
\end{proposition}

\begin{proof}
\leanok
This follows from Proposition \ref{prop.fps.binary-product-rule} by
observing that both formulations sum over the same collection of finite
subsets of $\mathbb{N}$.
\end{proof}

\subsection{Another example}

\begin{proposition}
[Euler]\label{prop.gf.prod.euler-odd}
\lean{euler_odd_parts_identity}
\leantarget
\leanok
We have%
\[
\prod_{i>0}\left(  1-x^{2i-1}\right)  ^{-1}=\prod_{k>0}\left(  1+x^{k}\right)
.
\]
\end{proposition}

\begin{proof}
For each $k>0$, we have $1+x^{k}=\dfrac{1-x^{2k}}{1-x^{k}}$ (since
$1-x^{2k}=\left(  1-x^{k}\right)  \left(  1+x^{k}\right)  $). Thus,%
\begin{align*}
\prod_{k>0}\left(  1+x^{k}\right)
&  =\prod_{k>0}\dfrac{1-x^{2k}}{1-x^{k}}
=\dfrac{\prod_{k>0}\left(  1-x^{2k}\right)  }{\prod_{k>0}\left(
1-x^{k}\right)  }\\
&  =\dfrac{1}{\left(  1-x^{1}\right)  \left(  1-x^{3}\right)  \left(
1-x^{5}\right)  \left(  1-x^{7}\right)  \cdots}\\
&  =\prod_{i>0}\left(  1-x^{2i-1}\right)^{-1},
\end{align*}
where the even factors cancel. The cancellation is justified by the fact that
the family $\left(  1-x^{k}\right)  _{k>0}$ factors as a product over even and odd indices,
by the partition of positive integers into evens and odds
(by Proposition \ref{prop.fps.union-mulable}).

\textbf{Lean proof strategy.}
The Lean proof proceeds differently: both sides are shown to equal the same
partition-theoretic generating function.
The RHS equals $\sum_n d_n \cdot x^n$
(generating function for partitions into distinct parts, i.e., parts of multiplicity $<2$).
The LHS equals $\sum_n o_n \cdot x^n$
(generating function for partitions into odd parts).
These are equal by Glaisher's theorem.
\end{proof}

\begin{theorem}
[Euler]\label{thm.gf.prod.euler-comb}
\lean{euler_distinct_odd_partitions}
\leantarget
\leanok
Let $n\in\mathbb{N}$. Then, $d_{n}=o_{n}%
$, where

\begin{itemize}
\item $d_{n}$ is the number of ways to write $n$ as a sum of distinct positive
integers, without regard to the order;

\item $o_{n}$ is the number of ways to write $n$ as a sum of (not necessarily
distinct) odd positive integers, without regard to the order.
\end{itemize}
\end{theorem}

\begin{proof}
Proposition \ref{prop.gf.prod.euler-odd} shows that
the generating functions for $d_n$ and $o_n$ are equal, hence $d_n = o_n$ for all $n$.
\end{proof}

\begin{proposition}
\label{prop.gf.prod.euler-odd-alt}
\lean{euler_odd_parts_identity'}
\leanhelper
\leanok
Alternative formulation of Euler's identity: indexing the left-hand side
over odd positive integers rather than all positive integers,
\[
\prod_{\substack{i\in\mathbb{N}^{+}\\i\text{ odd}}}
\left(1-x^i\right)^{-1} = \prod_{k>0}\left(1+x^k\right).
\]
\end{proposition}

\begin{proof}
\leanok
This follows from Proposition \ref{prop.gf.prod.euler-odd} by reindexing
the product using the bijection $i\mapsto 2i-1$ between $\mathbb{N}^{+}$
and the odd positive integers.
\end{proof}

\subsection{Infinite products and substitution}

\begin{lemma}
\label{lem.fps.continuous-rescale}
\lean{continuous_rescale}
\leanhelper
\leanok
For any $a\in K$, the rescaling map $\mathrm{rescale}(a) : K[[x]]\to K[[x]]$
defined by $f(x)\mapsto f(ax)$ is continuous (with respect to the
coefficient-wise topology on $K[[x]]$).
\end{lemma}

\begin{proof}
The $n$-th coefficient of $\mathrm{rescale}(a)(f)$ is $a^n\cdot [x^n]f$.
This is continuous because $[x^n]$ is continuous (it is a projection)
and scalar multiplication by $a^n$ is continuous.
\end{proof}

\begin{proposition}
\label{prop.fps.subs.rule-infprod}
\lean{tprod_rescale_substitution}
\leantarget
\leanok
If $\left(  f_{i}\right)  _{i\in I}\in
K\left[  \left[  x\right]  \right]  ^{I}$ is a multipliable family of FPSs,
and if $g\in K\left[  \left[  x\right]  \right]  $ is an FPS satisfying
$\left[  x^{0}\right]  g=0$, then the family $\left(  f_{i}\circ g\right)
_{i\in I}\in K\left[  \left[  x\right]  \right]  ^{I}$ is multipliable as well
and we have $\left(  \prod_{i\in I}f_{i}\right)  \circ g=\prod_{i\in I}\left(
f_{i}\circ g\right)  $.
\end{proposition}

\begin{proof}
The Lean formalization proves the special case of rescaling $g = ax$
(i.e., substituting $ax$ for $x$).
The proof uses Lemma \ref{lem.fps.continuous-rescale} (the rescale map is
continuous) together with the fact that it is a ring homomorphism,
so it commutes with infinite products.
\end{proof}

\subsection{Exponentials, logarithms and infinite products}

\begin{proposition}
\label{prop.fps.Exp-Log-infsum}
\lean{PowerSeries.Exp_sum, PowerSeries.Exp_multipliable_of_summable}
\leantarget
\leanok
Assume that $K$ is a commutative $\mathbb{Q}%
$-algebra. Let $\left(  f_{i}\right)  _{i\in I}\in K\left[  \left[  x\right]
\right]  ^{I}$ be a summable family of FPSs in $K\left[  \left[  x\right]
\right]  _{0}$. Then, $\left(  \operatorname{Exp}f_{i}\right)  _{i\in I}$ is
a multipliable family of FPS in $K\left[  \left[  x\right]  \right]  _{1}$,
and we have $\sum_{i\in I}f_{i}\in K\left[  \left[  x\right]  \right]  _{0}$
and%
\[
\operatorname{Exp}\left(  \sum_{i\in I}f_{i}\right)  =\prod_{i\in I}\left(
\operatorname{Exp}f_{i}\right)  .
\]
\end{proposition}

\begin{proof}
The $\operatorname{Exp}$ map converts infinite sums in $K[[x]]_0$
to infinite products in $K[[x]]_1$. This is because $\operatorname{Exp}$
is a continuous ring homomorphism from $K[[x]]_0$ (with addition) to
$K[[x]]_1$ (with multiplication), and therefore commutes with infinite sums
and products.
\end{proof}

\begin{proposition}
\label{prop.fps.Exp-Log-infprod}
\lean{PowerSeries.Log_tprod, PowerSeries.Log_summable_of_multipliable}
\leantarget
\leanok
Assume that $K$ is a commutative $\mathbb{Q}%
$-algebra. Let $\left(  f_{i}\right)  _{i\in I}\in K\left[  \left[  x\right]
\right]  ^{I}$ be a multipliable family of FPSs in $K\left[  \left[  x\right]
\right]  _{1}$. Then, $\left(  \operatorname{Log}f_{i}\right)  _{i\in I}$ is
a summable family of FPS in $K\left[  \left[  x\right]  \right]  _{0}$, and we
have $\prod\limits_{i\in I}f_{i}\in K\left[  \left[  x\right]  \right]  _{1}$
and%
\[
\operatorname{Log}\left(  \prod\limits_{i\in I}f_{i}\right)  =\sum_{i\in
I}\left(  \operatorname{Log}f_{i}\right)  .
\]
\end{proposition}

\begin{proof}
The $\operatorname{Log}$ map converts infinite products in $K[[x]]_1$
to infinite sums in $K[[x]]_0$. This follows from the fact that
$\operatorname{Log}$ is a continuous ring homomorphism from $K[[x]]_1$
(with multiplication) to $K[[x]]_0$ (with addition), and therefore commutes
with infinite products and sums.
\end{proof}
